\chapter{Functoriality, restriction, and trace}
\index{functoriality!comparison theories}\index{restriction!psh function}\index{trace!comparison}
\label{chap:functoriality-trace}

The comparison dictionary is useful only when it respects geometric
maps. Birational pullback is usually formal, finite pushforward carries
a degree, restriction may destroy singular metrics, and analytic trace
is not literal restriction. This chapter separates these operations and
proves the two geometric cases in which the trace Dirichlet problem is
currently controlled.

There are four comparison maps to keep distinct. Pullback acts on
potentials and currents; pushforward acts on measures and carries the
generic degree; a determinant norm sends line bundles and their metrics in
the opposite direction; and analytic trace is a limiting operation on
singularity types. Only the last map has target
$\QPSH(\widetilde Y)/{\sim_P}$, and its representatives are
$\mathcal I$-good. In particular, neither a
set-theoretic restriction nor a determinant norm can be substituted for it.
The point of the chapter is to make the input and output of each arrow
visible before any formula is used.

Unless stated otherwise, complex spaces in this chapter are connected
compact K\"ahler manifolds. A finite-map statement in the algebraic
category assumes normal projective varieties. For a nonempty connected closed submanifold
$Y\subset X$, write
\[
 D_Y:=\{\varphi\in\QPSH(X):\nu(\varphi,Y)=0\}.
\tag{67.1}\label{eq:functor-trace-domain}
\]
This is the natural domain of the analytic trace constructed in
\Cref{chap:operations-trace}.
All Monge--Amp\`ere measures in this chapter have raw intersection mass.
For an ordinary complex norm weight $h=-\log\|e\|$, the curvature is
$\ddc(2h)$ in the book's convention. Probability measures are obtained
only by dividing each measure by its own polarization volume.

\section{Birational pullback}
\label{sec:functor-birational}

Let $\pi:X'\to X$ be a proper bimeromorphic morphism between smooth
compact K\"ahler manifolds. If
\[
                         T=\theta+\ddc\varphi\geq0,
\]
then
\[
             \pi^*T=\pi^*\theta+\ddc(\varphi\circ\pi)\geq0.
\tag{67.2}\label{eq:functor-current-pullback}
\]

\begin{proposition}\label{prop:functor-birational-pullback}
With the notation above, for $1\leq p\leq n=\dim X$:
\begin{enumerate}
\item generic Lelong numbers satisfy
      $\nu(\pi^*T,E)=\nu(T,v_E)$ for every prime divisor
      $E\subset X'$;
\item multiplier ideals satisfy
      \[
       \mathcal I(m\varphi)
       =
       \pi_*\!\left(
          \mathcal O_{X'}(K_{X'/X})
          \otimes\mathcal I(m\varphi\circ\pi)
       \right);
      \tag{67.3}\label{eq:functor-transformation-ideal}
      \]
\item nonpluripolar products satisfy
      \[
       \pi_*\left\langle(\pi^*T)^p\right\rangle
                   =\left\langle T^p\right\rangle;
      \tag{67.4}\label{eq:functor-npp-pullback}
      \]
\end{enumerate}
\par
\end{proposition}

\begin{proof}
For the first assertion, $v_E$ means $\operatorname{ord}_E$ without a
further positive scaling. The coefficient of $[E]$ in the Siu
decomposition of $\pi^*T$ is its generic Lelong number. By the valuative
definition in \Cref{prop:ch51-divisorial-calculus}, this is $\nu(T,v_E)$.
A local potential has that logarithmic coefficient in the divisorial
sense; its remaining term need not be bounded. Smooth changes of local
potential do not alter the number.

For the second assertion, use the multiplier-ideal convention
$|g|^2\rme^{-m\varphi}\in L^1_{\mathrm{loc}}$ of
\Cref{def:complex-multiplier-ideal}. On coordinate charts write $J_\pi$
for the holomorphic Jacobian. Change of variables away from the analytic
exceptional sets gives
\[
 \int_U |g|^2\rme^{-m\varphi}\,\rd V_X
 =\int_{\pi^{-1}(U)}
    |(g\circ\pi)J_\pi|^2\rme^{-m\varphi\circ\pi}\,\rd V_{X'},
\tag{67.4a}\label{eq:functor-jacobian-integrability}
\]
with the volume densities transported by the coordinate Jacobian;
comparably bounded smooth densities give the same integrability test.
The equality remains valid for nonnegative integrands and infinite
integrals. A local section of $\mathcal O(K_{X'/X})$ has meromorphic
coefficient $h/J_\pi$. It belongs to the twisted multiplier ideal exactly
when $h$ is holomorphic and $|h|^2\rme^{-m\varphi\circ\pi}$ is locally
integrable. The coefficients of a section patch as one meromorphic function on
$\pi^{-1}(U)$. Off the exceptional image it descends to a holomorphic
function on $U$. The image has codimension at least two, so normality
extends that function holomorphically across it. Conversely, a
holomorphic $g$ downstairs has coefficient $g\circ\pi$ upstairs, and
its twisted-ideal membership tests $(g\circ\pi)J_\pi$. Properness and a finite chart cover over a
compact neighbourhood let the integrability equivalence be checked in
both directions. This proves \eqref{eq:functor-transformation-ideal}
for arbitrary quasi-psh weights, without resolving their singularities.
For an analytic singularity, the same formula reduces to the usual
relative-canonical-divisor calculation on a log resolution.

For \eqref{eq:functor-npp-pullback}, choose an analytic exceptional pair
$Z\subset X$ and $F=\pi^{-1}(Z)$ such that
$\pi:X'\setminus F\to X\setminus Z$ is biholomorphic. On this open set,
biholomorphic invariance of the bounded local truncations gives equality
of the nonpluripolar products. Both products have zero trace mass on their
analytic exceptional sets. To check this after pushforward, for a smooth
Hermitian form $\omega$ on $X$ and a K\"ahler form $\omega'$ on $X'$,
compactness gives $\pi^*\omega\leq C\omega'$. Positivity and the smooth
projection formula show
\[
 \left(\pi_*\langle(\pi^*T)^p\rangle\wedge\omega^{n-p}\right)(Z)
 \leq C^{n-p}
 \left(\langle(\pi^*T)^p\rangle\wedge(\omega')^{n-p}\right)(F)=0.
\tag{67.4b}\label{eq:functor-npp-exceptional-trace}
\]
Equality off $Z$ therefore extends across $Z$. This is also the mixed
modification identity \eqref{eq:complex-npp-modification} proved in
\Cref{chap:complex-prerequisites}.

\end{proof}

The proposition deliberately contains no analytic strict-transform claim:
its data do not specify a trace center.  For a transcendental
non-Archimedean metric satisfying the trace-domain hypotheses, the sourced
birational identity is the separate statement
\eqref{eq:operations-na-trace-birational}.

\begin{example}[The blowup Jacobian]
\label{ex:functor-blowup-jacobian}
Let $X=\mathbf C^2$, let $\pi:X'\to X$ be the blowup of the origin, and
take
\[
                  \varphi=c\log(|x|^2+|y|^2)+O(1),
                  \qquad c>0.
\]
On the chart $y=xu$, one has
\[
 \varphi\circ\pi
   =c\log|x|^2+c\log(1+|u|^2)+O(1).
\]
The second term is locally bounded, so the exceptional divisor $E=\{x=0\}$
has generic Lelong number $c$. Since $K_{X'/X}=E$, the extra copy of
$E$ in \eqref{eq:functor-transformation-ideal} is precisely the Jacobian
factor $|x|^2$ of the blowup chart. This two-dimensional calculation is
the local model for both the valuation identity and the multiplier-ideal
transformation rule. More explicitly, put $a=mc$ and let
$k$ be the order of a holomorphic germ at the origin. Its radial
integrability condition is
\[
 \int_0^\epsilon r^{2k-2a+3}\,\rd r<\infty
 \quad\Longleftrightarrow\quad k>a-2.
\]
Hence
\[
 \mathcal I(m\varphi)=\mathfrak m_0^{\max\{\lfloor a\rfloor-1,0\}}.
\tag{67.4c}\label{eq:functor-blowup-multiplier-example}
\]
Upstairs the ideal is $\mathcal O(-\lfloor a\rfloor E)$; twisting by
$K_{X'/X}=E$ and pushing down gives the same power of $\mathfrak m_0$,
including the integer threshold. For $a<1$, the pushforward of
$\mathcal O(E)$ is $\mathcal O_X$ by normality. This is a local calculation;
compactness is not required for it.
\end{example}

\begin{corollary}\label{cor:functor-full-mass-birational}
If $\alpha$ is big and $T\in\alpha$, then
\[
 \int_{X'}\left\langle(\pi^*T)^n\right\rangle
       =
 \int_X\left\langle T^n\right\rangle.
\tag{67.5}\label{eq:functor-full-mass}
\]
Thus $T$ has full Monge--Amp\`ere mass if and only if $\pi^*T$
does.
\end{corollary}

\begin{proof}
Take $p=n$ in \eqref{eq:functor-npp-pullback} and integrate; the
projection formula gives \eqref{eq:functor-full-mass}.  A current with minimal singularities pulls back to one with minimal
singularities by \cite[Proposition~1.12]{BEGZ10}. Applying the same mass
identity to that current gives
$\vol(\pi^*\alpha)=\vol(\alpha)$.  Equality of the two masses therefore
shows that one current has full mass exactly when the other does.
\end{proof}

\begin{remark}[Descent in a pulled-back class]
\label{rem:functor-no-inverse-pullback}
A current in an arbitrary class on $X'$ need not descend; the exceptional
current $[E]$ on a blowup is an example. A positive current in a pulled-back
class is different. If $T'=\pi^*\theta+\ddc u\geq0$, then $u$ is constant
on every connected fibre, with $-\infty$ allowed. Indeed its local psh
potential restricts to a psh function on each compact fibre; the maximum
principle and connectedness give constancy. It descends to a $\theta$-psh
potential on $X$: off the exceptional image this is immediate, and local
upper bounds permit psh extension across that image. Thus $T'=\pi^*T$.
This is the connected-fibre assertion following
\cite[Proposition~1.12]{BEGZ10}. The pulled-back cohomology hypothesis
supplies the descent in this positive-current setting.
\end{remark}

On the non-Archimedean side, a dominating model simply refines the
same valuation space. If $f:Y\to X$ is a morphism of projective
varieties, composition of semivaluations gives
\[
                f^{\An}:Y^{\An}\longrightarrow X^{\An}.
\tag{67.6}\label{eq:functor-analytic-map}
\]
A model metric pulls back through the pulled-back model line bundle.
The next section states the measure comparison with its dimension and
generic degree fixed.

\section{Finite maps, degree, and norms}
\label{sec:functor-finite}

Let $f:Y\to X$ be finite surjective between normal projective
varieties, of generic degree
\[
                         d=[K(Y):K(X)].
\tag{67.7}\label{eq:functor-degree}
\]
No separability hypothesis is imposed.

\begin{proposition}\label{prop:functor-finite-ma}
Let $L$ be ample on $X$, and let $\phi$ be a continuous
semipositive metric on $L$, either complex or non-Archimedean. In the
complex case, assume in addition that $X$ and $Y$ are projective
manifolds. Then
\[
        f_*\operatorname{MA}(f^*\phi)
                   =d\,\operatorname{MA}(\phi).
\tag{67.8}\label{eq:functor-finite-ma}
\]
In particular,
\[
                  ((f^*L)^n)=d(L^n).
\tag{67.9}\label{eq:functor-finite-volume}
\]
\end{proposition}

\begin{proof}
For a smooth complex metric and a continuous test function $a$ on $X$,
the definition of pushforward gives
\[
 \int_X a\,\rd f_*\operatorname{MA}(f^*\phi)
   =\int_Y(a\circ f)(f^*c_1(L,\phi))^n.
\]
Outside the branch locus, $f$ is a covering of degree $d$, while the
branch locus and its inverse image have measure zero for the smooth top
wedge. The change-of-variables formula therefore makes the last integral
equal to $d\int_Xa\,c_1(L,\phi)^n$. This proves the measure identity;
taking $a=1$ gives \eqref{eq:functor-finite-volume}.

For the continuous complex case the argument is local.  After trivializing
$L$ on a coordinate ball, approximate its continuous psh weight from
above on smaller balls by smooth psh weights.  Their pullbacks give the
corresponding approximants on $Y$, so the smooth identity just proved
passes to the limit by Bedford--Taylor monotone continuity.  This avoids
assuming a global uniform approximation by smooth semipositive metrics.

For a semipositive model metric, the mass at a divisorial point is an
intersection number on the corresponding component of the special fibre.
The finite algebraic projection formula, with residue degrees and
ramification multiplicities included, gives the factor $d$ after summing
over the components of the source special fibre lying above a given
component of the target; one may work on the normalized graph of the model morphism.
Its pushforward of the special-fibre cycle is $d$ times the target
special-fibre cycle. Intersecting with the pulled-back model line bundle
and applying the projection formula gives the stated coefficient identity.
This uses the whole cycle, including its multiplicities. Continuous non-Archimedean semipositive metrics are uniform
limits of semipositive model metrics, and their Monge--Amp\`ere measures vary
weakly continuously under that approximation. Testing against $a$ and
passing to the limit proves \eqref{eq:functor-finite-ma} in the
non-Archimedean setting as well.
\end{proof}

\begin{example}[A branched cover of the projective line]
\label{ex:functor-branched-ma}
Let $f:\mathbf P^1\to\mathbf P^1$ be $z\mapsto z^d$, $d\geq1$, and
let $\omega$ be the Fubini--Study form of mass one. On the affine chart,
\[
 f^*\omega=
 \frac{d^2|z|^{2d-2}}{\pi(1+|z|^{2d})^2}\,\rd A(z).
\tag{67.9a}\label{eq:functor-branched-density}
\]
For a continuous test $a$, polar coordinates and $u=\rho^d$ give
\[
 \int a(z^d)f^*\omega
 =\frac d\pi\int_0^\infty\int_0^{2\pi}
       a(u\rme^{i\vartheta})\frac{u}{(1+u^2)^2}
       \,\rd\vartheta\,\rd u=d\int a\,\omega.
\]
The angular substitution counts all $d$ turns. In particular the raw
source mass is $d$, whereas the probability measure $d^{-1}f^*\omega$
pushes to $\omega$. Branch points add no atomic mass.
\end{example}

\begin{remark}[Singular metrics]
\label{rem:functor-finite-singular}
For singular finite-energy metrics, the same formula requires a
specified nonpluripolar extension and an approximation theorem.
Equation~\eqref{eq:functor-finite-ma} is not being asserted for an
arbitrary extended-real psh function with no mass hypothesis.
\end{remark}

The reverse direction uses the norm, not an average chosen pointwise.
Suppose first that $f$ is finite locally free of rank $d$. For a
line bundle $M$ on $Y$, define
\[
 \operatorname{Nm}_f(M)
   :=\det(f_*M)\otimes\det(f_*\mathcal O_Y)^{-1}.
\tag{67.10}\label{eq:functor-norm-line}
\]
For a local section $s$ of $M$, multiplication by $s$ induces
the norm section $\operatorname{Nm}_f(s)$.

\begin{proposition}\label{prop:functor-determinant-norm}
Let $f$ be finite locally free of rank $d$, and let
$\|\cdot\|_M$ be a continuous metric on $M$. The determinant norm
defines a continuous metric on $\operatorname{Nm}_f(M)$. At an
analytic point $x\in X^{\An}$,
\[
 -\log\|\operatorname{Nm}_f(s)(x)\|
   =
 \sum_{y\in f^{-1}(x)}
       \ell_y\,\bigl(-\log\|s(y)\|_M\bigr),
\tag{67.11}\label{eq:functor-norm-weight}
\]
where
\[
   \ell_y=\dim_{\mathcal H(x)}
          \mathcal O_{Y_x,y}
\tag{67.12}\label{eq:functor-fibre-length}
\]
is the dimension of the local fibre algebra over the completed base
residue field. If its residue field is $K_y$ and its Artin length is $m_y$,
then $\ell_y=m_y[K_y:\mathcal H(x)]$.
\end{proposition}

\begin{proof}
Locally on the base, choose a frame $e$ of $M$ over the whole inverse
image. Such frames exist after shrinking around a fibre: an invertible
module over the finite semilocal fibre algebra is free, and Nakayama's
lemma lifts its generator. Define the weight of the frame $\operatorname{Nm}_f(e)$ by
\[
 H_e(x)=\sum_{y\in f^{-1}(x)}\ell_y h_e(y),
 \qquad h_e(y)=-\log\|e(y)\|_M.
\tag{67.12a}\label{eq:functor-norm-frame-weight}
\]
For an invertible function $b$, the algebraic determinant satisfies
\[
 |\operatorname{Nm}(b)(x)|
   =\prod_y |b(y)|^{\ell_y}.
\tag{67.12b}\label{eq:functor-norm-absolute-value}
\]
To see the residue-degree factor, filter the local Artin algebra by
powers of its maximal ideal. Its determinant is
$N_{K_y/\mathcal H(x)}(\bar b)^{m_y}$; the absolute value of the field
norm is $|b(y)|^{[K_y:\mathcal H(x)]}$.
Thus the frame-change formula for $H_e$ is exactly the norm-line
transition formula. The same determinant calculation for a section
$s=be$, with zero values allowed, proves
\eqref{eq:functor-norm-weight} in the extended-real sense.

Here continuity requires more than polynomiality of the determinant:
$h_e$ is an arbitrary continuous weight. Fix a fibre
$f^{-1}(x)=\{y_1,\ldots,y_q\}$ and $\epsilon>0$. Choose disjoint
neighbourhoods $W_i$ of $y_i$ on which
$|h_e-h_e(y_i)|<\epsilon$. Properness permits shrinking the base so
that every nearby fibre lies in $\bigcup_iW_i$. The total fibre
multiplicity in $W_i$ is then $\ell_{y_i}$. Here is the finite-flat
local argument for that constancy. The finite fibre algebra splits into
local factors. Functions regular near $x$, with denominators nonzero
at $x$, are dense in its completed residue field. Approximate the
factor idempotents by elements of the analytic
algebra on a base neighbourhood. Their multiplication matrices have
characteristic polynomials varying continuously with the base point.
The roots near $1$ and $0$ stay in disjoint clusters, with the same
multiplicities. Over a non-Archimedean completed residue field this is
monic-polynomial factorization across disjoint spectral discs; over
$\mathbf C$ it is continuity of roots counted with multiplicity.
The ranks of these clusters count exactly the base-field dimensions
$\ell_{y_i}$, including nilpotents and residue degrees. Consequently
\[
 |H_e(x')-H_e(x)|\leq d\epsilon
\tag{67.12c}\label{eq:functor-norm-continuity-bound}
\]
for all sufficiently near $x'$. This proves continuity of the local
norm weight and hence of the metric.
\end{proof}

\begin{example}[A radicial fibre]
\label{ex:functor-radicial-norm}
In characteristic $p>0$, take $f:\mathbf A^1\to\mathbf A^1$ given by
$x=y^p$. Over a point with residue value $a$, the fibre algebra is
$k[y]/(y^p-a)$. If $a$ is not a $p$th power in $k$, this is a field with Artin
length one and residue degree $p$; after adjoining a $p$th root it is
$k'[\epsilon]/(\epsilon^p)$, of Artin length $p$ and residue degree one.
In both cases the dimension $\ell_y$ is $p$. Hence a
constant fibre value $c$ contributes $p(-\log|c|)$ to
\eqref{eq:functor-norm-weight}. Counting the unique point with weight one
would lose the entire generic degree.
\end{example}

\begin{example}[Continuity when two points coalesce]
\label{ex:functor-coalescing-norm}
On the complex affine double cover $x=y^2$, give the trivial line its
continuous frame weight $h(y)=a+|y|^2$. For $x\neq0$, the two roots have
weights $a+|x|$, while the fibre at zero is
$\mathbf C[y]/(y^2)$ of length two. Thus the norm frame has the continuous
weight $H(x)=2a+2|x|$ on the whole chart. The norm of the section $y$ is
$-x$, so its logarithmic norm weight is $-\log|x|+2a+2|x|$, with value
$+\infty$ at zero. Continuity of the metric concerns its nonvanishing
frame; a vanishing section retains its logarithmic pole.
\end{example}

\begin{remark}[Inseparability and nilpotents]
\label{rem:functor-inseparable}
Formula \eqref{eq:functor-norm-weight} sums lengths, not distinct
geometric points. It therefore retains purely inseparable
multiplicity and nonreduced fibre structure. Replacing $\ell_y$ by
$1$ would silently assume reduced separable fibres.
\end{remark}

For a finite map which is not flat, flattening the finite algebra
produces a birational square of the form
\[
\begin{tikzcd}[column sep=large,row sep=large]
W \arrow[r,"\alpha"] \arrow[d,"q"']&
Y \arrow[d,"f"]\\
Z \arrow[r,"\beta"']&
X ,
\end{tikzcd}
\tag{67.13}\label{eq:functor-finite-flat-square}
\]
where $\beta$ is a projective modification, $W$ is the strict
transform of $Y$, and $q$ is finite locally free. The strict transform is
the schematic closure of the inverse image over the dense flat locus;
it need not be the whole fibre product. These are the exact hypotheses
and construction of \Cref{lem:bj-finite-flatification}. The determinant norm is formed
upstairs and then compared birationally. This is the mechanism behind
the finite descent theorem of
\Cref{chap:bj-finite-descent}; it is not a claim that $f$
itself becomes flat without modifying $X$.

\section{Base change and descent}
\label{sec:functor-base-change}

Consider the ramified base change
\[
                         t=s^m.
\tag{67.14}\label{eq:functor-ramified-base-change}
\]
If $|t|=r$, the compatible norm on the new base has
$|s|=r^{1/m}$. A valuation normalized by $v(t)=1$ pulls back to a
valuation with $v(s)=m^{-1}$.

\begin{proposition}\label{prop:functor-scaling-ledger}
Let a normal test configuration be pulled back by
\eqref{eq:functor-ramified-base-change} and then normalized.
\begin{enumerate}
\item With the original ray parameter $s_t=-\log|t|$, the associated
      ray is unchanged.
\item With the new parameter $s_s=-\log|s|$, its speed is multiplied
      by $m$.
\item Non-Archimedean functionals homogeneous of degree one under
      rescaling of the ray parameter satisfy
      \[
             F^{\NA}(\phi^{\NA}_{(m)})
                    =mF^{\NA}(\phi^{\NA})
      \tag{67.15}\label{eq:functor-scaling-functional}
      \]
              when the left-hand side is read in the new normalization
        $v(s)=1$ and the right-hand side in the old one $v(t)=1$ ---
        that is, when each test configuration is read with its own base
        coordinate normalized to one.  In one \emph{fixed} normalization the
        two sides agree, which is item~(1) restated: expressing both
        arguments with $v(s)=1$ multiplies both by $m$, and a
        degree-one homogeneous functional then preserves the equality.  The
        factor $m$ belongs to the change of normalization, not to either
        metric.
\end{enumerate}
\par
\end{proposition}

\begin{proof}
The identity $s_t=ms_s$ is immediate from $t=s^m$.  The same
relation gives $\ord_E(t)=m\ord_E(s)$ for every component $E$ of the
special fibre, so the multiplicity of $E$ in the $v(t)=1$ convention is
$m$ times its multiplicity in the $v(s)=1$ convention.  A model
function is a vertical coefficient divided by that multiplicity.  Its
$v(t)=1$ value is therefore unchanged by the base change, which is
item~(1), while its $v(s)=1$ value is multiplied by $m$, which is
item~(2).  Applying this to a functional homogeneous of degree one in the
ray parameter gives \eqref{eq:functor-scaling-functional}.  The factor
$m$ comes from the projection formula for the degree-$m$ map of total
spaces; it is not a consequence of degree-one homogeneity of the
intersection numbers themselves, which for $E^{\NA}$ are of degree
$n+1$.
\end{proof}

\begin{example}[Base change scales the valuation argument]
\label{ex:functor-base-change-action}
On $X^{\mathrm{an,triv}}$, write $\widehat\phi$ for a unit-normalized
metric relative to the trivial reference. Its degree-$m$ base-change
scaling is
\[
 D_m\widehat\phi(v)=m\widehat\phi(m^{-1}v).
\tag{67.15a}\label{eq:functor-base-change-action}
\]
This is the action of \cite[Section~6.3]{BHJ17}; it scales the valuation
argument as well as the value. For the two-section metric on
$\mathcal O_{\mathbf P^1}(1)$, in the affine region $v(z)\geq0$ the
relative weight is $\widehat\phi(v)=\max\{0,1-v(z)\}$. There
\[
 D_m\widehat\phi(v)=\max\{0,m-v(z)\}.
\]
At $v(z)=1/2$ and $m=2$, its value is $3/2$, whereas
$2\widehat\phi(v)=1$. Thus degree-one homogeneity in
\eqref{eq:functor-scaling-functional} refers to $D_m$, not pointwise
scalar multiplication of a fixed function on the valuation space.
The ray runs at speed $m$ with the new radial parameter, while the
valued-field pullback is unchanged when $|s|=r^{1/m}$ is used.
\end{example}

\begin{remark}[Descent datum]
\label{rem:functor-descent-datum}
A metric after finite base change does not automatically descend.
One needs invariance under the descent relation, or under the Galois
group when a Galois presentation is available. In positive
characteristic, a purely inseparable factor has no useful Galois
group; descent must be formulated through the finite algebra or the
universal homeomorphism.
\end{remark}

\begin{proposition}\label{prop:functor-universal-homeomorphism}
If $f:Y\to X$ is finite, radicial, and surjective, then
\[
                    f^{\An}:Y^{\An}\longrightarrow X^{\An}
\tag{67.16}\label{eq:functor-universal-homeomorphism}
\]
is a homeomorphism. It need not identify line bundles or their model
structures without a descent datum.
\end{proposition}

\begin{proof}
Every support prime has a unique lift. The induced residue extension
is purely inseparable, and an absolute value has a unique extension. In
positive characteristic it is given, for a sufficiently large $e$, by
\[
                         |a|=|a^{p^e}|^{1/p^e}
\]
after $a^{p^e}$ has descended to the base residue field; in
characteristic zero the residue extension is trivial, so there is
nothing to extend. Thus \eqref{eq:functor-universal-homeomorphism}
is a continuous bijection. Properness makes the source compact and
the target Hausdorff, hence it is a homeomorphism. The last statement
is algebraic: a homeomorphism of analytifications contains no
linearization or cocycle for a line bundle.
\end{proof}

\section{Restriction is not trace}
\label{sec:functor-restriction}

If $\varphi$ is quasi-psh on $X$, its literal restriction to
$Y$ exists only when $\varphi|_Y\not\equiv-\infty$. Even then it
depends on the chosen representative, not merely its singularity
class. The analytic trace instead restricts a quasi-equisingular
approximation and takes a strong limit.

The codomain is part of the theorem. Multiplier ideals control the input:
trace depends only on the $\mathcal I$-class of the ambient potential.  The
restriction theorem then gives the finer output: trace representatives are
$\mathcal I$-good and unique up to $P$-equivalence. Thus the quotient map
\[
 \{P\text{-classes in }D_Y\}\longrightarrow
 \{\mathcal I\text{-classes in }D_Y\}
\]
may be applied before trace, while the target lies in the $P$-quotient.

\begin{definition}[Analytic trace]
\label{def:functor-analytic-trace}
For $\varphi\in D_Y$, choose a quasi-equisingular approximation
$(\varphi_j)$ with analytic singularities. On the normalization
$\nu:\widetilde Y\to Y$, define
\[
                   \Tr_Y(\varphi)
       :=d_S\text{-}\lim_{j\to\infty}
                   \nu^*(\varphi_j|_Y),
\tag{67.17}\label{eq:functor-analytic-trace}
\]
after fixing one sufficiently positive background class on
$\widetilde Y$, as required for the positive-mass $d_S$-space,
where equality of outputs is understood up to
$P$-equivalence.  Every output is $\mathcal I$-good.
\end{definition}

\begin{remark}[Restricted ambient ideals are image ideals]
\label{rem:functor-restricted-image-ideal}
If an analytic approximant is represented by an ambient coherent ideal
$\mathfrak a\subset\mathcal O_X$, the ideal used after restriction is
\[
 \mathfrak a\mathcal O_Y
 :=\operatorname{im}\!\left(
   \mathfrak a\otimes_{\mathcal O_X}\mathcal O_Y
   \longrightarrow\mathcal O_Y\right).
\]
It is not the abstract tensor product, and it is not automatically the
intrinsic multiplier ideal of a naively restricted singularity.  The
restriction theorem supplies, when literal restriction is defined, only
the one-sided inclusion
\[
 \mathcal I(m\varphi|_Y)
 \subseteq \mathcal I(m\varphi)\mathcal O_Y;
\]
no equality is used in the trace construction.  This is why the
$d_S$-limit and the sourced uniqueness theorem, rather than an
identification of restricted ideals, determine the output $P$-class.
\end{remark}

\begin{proposition}\label{prop:functor-trace-restriction}
Let $Y\subset X$ be a smooth submanifold and
$\varphi\in D_Y$.
\begin{enumerate}
\item If $\varphi$ has analytic singularities with rational coefficients and
      $\varphi|_Y\not\equiv-\infty$, then
      \[
                \Tr_Y(\varphi)\sim_P\varphi|_Y.
      \tag{67.18}\label{eq:functor-trace-analytic}
      \]
\item For arbitrary $\varphi\in D_Y$, literal restriction can be identically $-\infty$ even though
      each analytic approximant restricts nontrivially, whereas
      the trace is well-defined by
      \Cref{thm:operations-trace-well-defined}.
\end{enumerate}
\par
\end{proposition}

\begin{proof}
For these rational analytic singularities, take the constant quasi-equisingular
approximation $\varphi_j=\varphi$. It has the required curvature bound
and satisfies the multiplier-ideal inclusions, since for
$\lambda'>\lambda$ one has
$\mathcal I(\lambda'\varphi)\subseteq\mathcal I(\lambda\varphi)$.
Every restriction equals $\varphi|_Y$, so its $d_S$-distance from that
potential is zero. Independence of approximation from
\Cref{thm:operations-trace-well-defined} proves
\eqref{eq:functor-trace-analytic}. This is also
\cite[Example~3.7]{DXTraceV2}. The rationality hypothesis matches the
analytic approximants in that source; the broader real-coefficient
terminology of \Cref{chap:complex-prerequisites} alone does not supply
this particular trace interface.

For the second assertion, \Cref{ex:functor-zero-lelong-pole} supplies the
pointwise obstruction. The trace theorem makes two choices of
quasi-equisingular approximation comparable by diagonal restriction
estimates, so their limits have the same $P$-class. A further distinction
concerns increasing families: the increasing family in
\Cref{rem:operations-trace-not-increasing} has a limiting trace mass
strictly smaller than the trace mass of its ambient limit. That family
is not a decreasing quasi-equisingular approximation. It shows the
failure of a general increasing-continuity rule. The pointwise
obstruction and the increasing mass gap are distinct phenomena.
\end{proof}

The distinction can be summarized by the diagram
\[
\begin{tikzcd}[column sep=huge,row sep=large]
D_Y
 \arrow[r,"\text{quasi-equisingular approximation}"]
    \arrow[ddr,dashed,"\text{literal restriction}"']
&
\{\text{analytic singularities}\}^{\mathbf N}
 \arrow[d,"\text{restrict termwise}"]\\
&
\QPSH(\widetilde Y)^{\mathbf N}
 \arrow[d,"\text{$d_S$-limit}"]\\
&
\QPSH(\widetilde Y)/{\sim_P}.
\end{tikzcd}
\tag{67.19}\label{eq:functor-trace-diagram}
\]
The dashed arrow is unavailable in general.

\begin{example}[A zero-Lelong pole]
\label{ex:functor-zero-lelong-pole}
Near a smooth hypersurface $Y=\{z=0\}$, the function
\[
                  \varphi=-\log\bigl(-\log|z|^2\bigr)
\]
has generic Lelong number zero but equals $-\infty$ on $Y$. Its generic Lelong condition is the local condition defining
$D_Y$, although literal restriction is meaningless. The
calculation separates the domain condition in
\eqref{eq:functor-trace-domain} from pointwise finiteness. With
$q=\log|z|^2<0$, the radial profile $-\log(-q)$ has derivatives
$-1/q>0$ and $1/q^2>0$, so it is psh. Its quotient by $q$ tends to zero,
proving the stated Lelong number. Every multiplier weight is locally
integrable, since substitution $u=-\log\rho^2$ gives
\[
 \int_0^\epsilon \rho(-\log\rho^2)^b\,\rd\rho
   =\tfrac12\int_{-\log\epsilon^2}^{\infty}u^b\rme^{-u}\,\rd u<\infty
 \quad(b>0).
\]
The bounded truncations $\max\{\varphi,-j\}$ decrease to this germ and
all have the trivial analytic singularity type. Their restrictions are
constants $-j$, whose $d_S$-class is the bounded class despite their
pointwise limit $-\infty$.

For a compact example, let $\eta$ be the explicit global
$\omega_{\mathrm{FS}}$-psh log--log potential in
\Cref{rem:ch51-loglog-explicit}, and put $X=\mathbf P^1\times Y_0$
for a connected compact K\"ahler manifold $Y_0$. The potential
$p_1^*\eta$ has generic Lelong number zero along
$Y=\{0\}\times Y_0$ and equals $-\infty$ there. Its global bounded
truncations $\max\{p_1^*\eta,-j\}$ have trivial multiplier ideals and
form a quasi-equisingular approximation in one product background.
They restrict to the constants $-j$. Thus the trace is the bounded
$P$-class on $Y_0$, whereas literal restriction is identically
$-\infty$.
\end{example}

\section{The split trace Dirichlet theorem}
\label{sec:functor-split-dirichlet}

The general trace Dirichlet problem asks whether a prescribed K\"ahler
current on $Y$ is the analytic trace of a K\"ahler current in a fixed
class on $X$. The next theorem is the exact product case validated
in the research ledger.

\begin{theorem}\label{thm:functor-split-trace}
Let $Y$ and $Z$ be connected compact K\"ahler manifolds, set
$X=Y\times Z$, and fix
\[
              \alpha=p_Y^*\alpha_Y+p_Z^*\alpha_Z
\tag{67.20}\label{eq:functor-split-class}
\]
with $\alpha_Y,\alpha_Z$ K\"ahler classes. Let
\[
             S=\theta_Y+\ddc\psi\in\alpha_Y
\tag{67.21}\label{eq:functor-prescribed-current}
\]
be a K\"ahler current, and let
$\eta\in\alpha_Z$ be a smooth K\"ahler form. Then
\[
                    T:=p_Y^*S+p_Z^*\eta
\tag{67.22}\label{eq:functor-product-current}
\]
is a K\"ahler current in $\alpha$, and for every $z_0\in Z$,
the analytic trace satisfies
\[
 \Tr_{Y\times\{z_0\}}(T)
       \sim_P P_{\alpha_Y}[\psi]_{\mathcal I}
       \sim_{\mathcal I}S.
\tag{67.23}\label{eq:functor-product-trace}
\]
If $S$ is $\mathcal I$-good, then its canonical
$\mathcal I$-model representative satisfies
\[
 \Tr_{Y\times\{z_0\}}(T)
 \sim_P P_{\alpha_Y}[\psi]_{\mathcal I}
 \sim_P S.
\tag{67.24}\label{eq:functor-product-p-good}
\]
\end{theorem}

\begin{proof}
\noindent\emph{Step 1: positivity in the fixed class.}
Choose $\varepsilon>0$ and a K\"ahler form $\omega_Y$ such that
\[
                         S\geq\varepsilon\omega_Y.
\]
Since $\eta$ is K\"ahler, it dominates
$\varepsilon'\omega_Z$ for some $\varepsilon'>0$. Therefore
\[
 T\geq
 \min\{\varepsilon,\varepsilon'\}
      (p_Y^*\omega_Y+p_Z^*\omega_Z),
\tag{67.25}\label{eq:functor-product-positive}
\]
so $T$ is a K\"ahler current.  Its class is $\alpha$ by
\eqref{eq:functor-split-class}.

\medskip
\noindent\emph{Step 2: pull back a quasi-equisingular approximation.}
Choose a decreasing quasi-equisingular approximation $(\psi_j)$ of
$\psi$ as in \eqref{eq:operations-trace-qe}. Since
$S\geq\varepsilon\omega_Y$, regularize in the background
$\theta_Y-\varepsilon\omega_Y$ and discard finitely many terms. Then
$\theta_Y+\ddc\psi_j\geq(\varepsilon/2)\omega_Y$ for all $j$.
The approximants have analytic singularities and converge pointwise to
$\psi$ with the required multiplier-ideal inclusions. Their strong
singularity limit is identified in Step~4; it need not be the original
$P$-class of $\psi$.
Set
\[
                          \Psi_j:=p_Y^*\psi_j.
\tag{67.26}\label{eq:functor-product-approximants}
\]
The projection $p_Y$ is smooth. Locally it is the projection
$U\times V\to U$; Fubini's theorem gives
\[
     \mathcal I(m\Psi_j)
                =p_Y^{-1}\mathcal I(m\psi_j)\cdot\mathcal O_X
\tag{67.27}\label{eq:functor-smooth-ideal-pullback}
\]
for every $m>0$. The same identity holds for $\psi$.
Consequently $(\Psi_j)$ is quasi-equisingular for $p_Y^*\psi$.
The centre $Y\times\{z_0\}$ has generic Lelong number zero: a generic
point of $Y$ lies outside the polar set of $\psi$. Thus $T$ belongs to
the trace domain before the limiting theorem is invoked.

\medskip
\noindent\emph{Step 3: restrict the approximants.}
Let $i_{z_0}:Y\hookrightarrow X$ be
$y\mapsto(y,z_0)$. Then
\[
                         i_{z_0}^*\Psi_j=\psi_j,
          \qquad i_{z_0}^*p_Z^*\eta=0.
\tag{67.28}\label{eq:functor-product-restriction}
\]
Since the smooth term $p_Z^*\eta$ changes no multiplier ideal, the
sequence used to compute the trace of $T$ restricts exactly to the
chosen sequence for $S$.

\medskip
\noindent\emph{Step 4: identify the limit.}
By the definition of analytic trace,
\[
       \Tr_{Y\times\{z_0\}}(T)
          \sim_P d_S\text{-}\lim_j\psi_j
          \sim_P P_{\alpha_Y}[\psi]_{\mathcal I}.
\tag{67.29}\label{eq:functor-product-limit}
\]
The uniform K\"ahler-current lower bound supplies a common positive
mass floor. The canonical-approximation theorem
\cite[Corollary~7.1.2]{XiaBook} therefore identifies
$d_S\text{-}\lim_j\psi_j$ with $P_{\alpha_Y}[\psi]_{\mathcal I}$.
Its multiplier ideals equal those of $\psi$, hence those of the
potential of $S$. This proves
\eqref{eq:functor-product-trace}. If $S$ is $\mathcal I$-good,
its $P$-class equals the $P$-class of its
$\mathcal I$-model envelope, giving the statement about the canonical
representative in
\eqref{eq:functor-product-p-good}.
\end{proof}

\begin{example}[Why the $Z$-term is necessary]
\label{ex:functor-product-vertical-positivity}
The current $p_Y^*S$ vanishes on tangent vectors to the $Z$-factor,
so it is not a K\"ahler current on $Y\times Z$ when $\dim Z>0$.
The term $p_Z^*\eta$ supplies precisely the missing vertical
positivity without changing the trace on a horizontal slice.
\end{example}

\begin{example}[Product mass with a logarithmic pole]
\label{ex:functor-product-mass}
Take $Y=Z=\mathbf P^1$ with Fubini--Study forms of mass one. For
$a\in\mathbf Q\cap(0,1)$, let $\psi=a\log\|s_p\|_{\mathrm{FS}}^2$, where $s_p$ is a
section of $\mathcal O(1)$ with divisor $p$. Then
\[
 S=\omega_Y+\ddc\psi=(1-a)\omega_Y+a\delta_p
       \geq(1-a)\omega_Y,
 \qquad T=p_Y^*S+p_Z^*\omega_Z.
\]
The raw class has self-intersection $2$, but the nonpluripolar product
removes the logarithmic pole on $\{p\}\times Z$. Locality and
multilinearity give
\[
 \langle T^2\rangle
   =2(1-a)p_Y^*\omega_Y\wedge p_Z^*\omega_Z,
 \qquad \int_X\langle T^2\rangle=2(1-a)>0.
\tag{67.29a}\label{eq:functor-product-npp-mass}
\]
On each horizontal slice, every pulled-back analytic approximant
restricts to the original approximant, and the transverse form vanishes.
Since $\psi$ has analytic singularities, the trace has its $P$-class by
\eqref{eq:functor-trace-analytic}. Strict positivity of $T$ and the trace
identity therefore coexist with a genuine loss of nonpluripolar mass.
\end{example}

\section{The retraction variant}
\label{sec:functor-retraction}

\begin{theorem}\label{thm:functor-retraction-trace}
Let $r:X\to Y$ be a holomorphic submersion of connected compact
K\"ahler manifolds with a holomorphic section
$i:Y\hookrightarrow X$, so $r\circ i=\operatorname{id}_Y$. Suppose
\[
                         \alpha=r^*\beta+\gamma
\tag{67.30}\label{eq:functor-retraction-class}
\]
where $\beta$ is K\"ahler and $\gamma$ has a smooth semipositive
representative $\eta$ satisfying
\[
 i^*\eta=0,
 \qquad
 \eta(\xi,J\xi)>0
       \quad\text{for }0\neq\xi\in\ker(\rd r).
\tag{67.31}\label{eq:functor-transverse-positive}
\]
If $S=\theta_Y+\ddc\psi\in\beta$ is a K\"ahler current, then
\[
                              T=r^*S+\eta
\tag{67.32}\label{eq:functor-retraction-current}
\]
is a K\"ahler current in $\alpha$, and
the analytic trace satisfies
\[
                       \Tr_{i(Y)}(T)
                 \sim_P P_\beta[\psi]_{\mathcal I}
                 \sim_{\mathcal I}S.
\tag{67.33}\label{eq:functor-retraction-trace}
\]
If $S$ is $\mathcal I$-good, then
$\Tr_{i(Y)}(T)\sim_P P_\beta[\psi]_{\mathcal I}\sim_P S$.
\end{theorem}

\begin{proof}
Choose a K\"ahler form $\omega_Y$ and $c>0$ with
$S\geq c\omega_Y$. First use the current inequality
$T\geq c\,r^*\omega_Y+\eta$. Its right-hand side is smooth.
For a nonzero tangent vector $\xi$, there are two cases. If $\rd r(\xi)\neq0$, then
$c\,r^*\omega_Y(\xi,J\xi)>0$. If
$\rd r(\xi)=0$, then $\eta(\xi,J\xi)>0$ by
\eqref{eq:functor-transverse-positive}. Since both summands are
semipositive, the smooth form $c\,r^*\omega_Y+\eta$ is pointwise positive definite. Compactness of the
unit tangent bundle supplies a uniform positive lower bound against any
fixed Hermitian metric. Hence $T$ is a K\"ahler current representing
$\alpha$.

Because $r$ is a holomorphic submersion, it is locally a smooth
projection. Thus multiplier ideals commute with pullback exactly as in
\eqref{eq:functor-smooth-ideal-pullback}. Pull a quasi-equisingular
approximation $(\psi_j)$ back by $r$. Since
$r\circ i=\operatorname{id}_Y$,
\[
                         i^*r^*\psi_j=\psi_j.
\tag{67.34}\label{eq:functor-retraction-approximation}
\]
The condition $i^*\eta=0$ says that the smooth transverse term
contributes no form on the section. The trace computation is therefore
identical to Steps $3$--$4$ in the proof of
\Cref{thm:functor-split-trace}. This gives
\eqref{eq:functor-retraction-trace} and the
$P$-equivalence conclusion for $\mathcal I$-good data.
\end{proof}

\begin{remark}[Why a bare retraction is insufficient]
\label{rem:functor-bare-retraction}
The identity $r\circ i=\operatorname{id}_Y$ provides an extension of the
singularity $\psi$, but it provides neither a decomposition
\eqref{eq:functor-retraction-class} nor positive curvature in
$\ker(\rd r)$. Thus a bare holomorphic retraction does not solve the
trace Dirichlet problem in a prescribed K\"ahler class.
\end{remark}

\section{Descent of trace and its exact sensitivity}
\label{sec:functor-trace-descent}

The comparison theorem from \Cref{chap:operations-trace} gives
\[
 \varphi\preceq_{\mathcal I}\psi
 \quad\Longrightarrow\quad
 \Tr_Y(\varphi)\preceq_P\Tr_Y(\psi)
\tag{67.35}\label{eq:functor-trace-order}
\]
on $D_Y$. Therefore the exact sourced map is
\[
\begin{tikzcd}[column sep=large]
D_Y/{\sim_{\mathcal I}}
 \arrow[r,"\Tr_Y"]
&
\QPSH(\widetilde Y)/{\sim_P}.
\end{tikzcd}
\tag{67.36}\label{eq:functor-trace-factorization}
\]

\begin{proposition}\label{prop:functor-domain-saturation}
The domain $D_Y$ is saturated for $\mathcal I$-equivalence.
Consequently,
\eqref{eq:functor-trace-factorization} is well-defined.
\end{proposition}

\begin{proof}
Equality of $\mathcal I$-types is equality of generic Lelong numbers
along divisorial valuations on modifications.  This is the equivalence
$\text{(i)}\Leftrightarrow\text{(ii)}$ of
\Cref{prop:va-i-equivalence}, which is stated there for a connected
projective manifold while the standing hypothesis here is compact
K\"ahler; the equivalence transfers because its proof uses only the
valuative characterization of multiplier ideals and strong openness, both
local. If $Y=X$, both generic Lelong numbers are zero. If $Y$ is a divisor,
use $\ord_Y$; in higher codimension, the exceptional divisor of the
blow-up of the smooth center $Y$ recovers $\nu(\varphi,Y)$. Hence
$\varphi\sim_{\mathcal I}\psi$ implies
$\nu(\varphi,Y)=\nu(\psi,Y)$. Applying
\eqref{eq:functor-trace-order} in both directions proves that the
arrow in \eqref{eq:functor-trace-factorization} is defined. The smooth-submanifold comparison is
\cite[Lemma~3.9]{DXTraceV2}; the background-independent $P$-class
interface is \cite[Propositions~8.1.1--8.1.2 and~8.2.1]{XiaBook}.
\end{proof}

The factorization can be read from left to right, but not backwards:
\[
 \frac{D_Y}{\sim_P}
 \longrightarrow
 \frac{D_Y}{\sim_{\mathcal I}}
 \xrightarrow{\ \Tr_Y\ }
 \frac{\QPSH(\widetilde Y)}{\sim_P}.
\]
The first arrow forgets information about the ambient input.  The second
arrow nevertheless has a $P$-valued target because restricted
quasi-equisingular approximants satisfy a finer uniqueness theorem.  This
does not recover the forgotten input $P$-class.

\begin{remark}[What the theorem does not recover]
\label{rem:functor-not-p-sensitive}
The source theorem does not assert that $\sim_{\mathcal I}$ and $\sim_P$
coincide on $D_Y$, nor does it assign an ordinary singularity type to the
trace.  Its conclusion is exactly the map in
\eqref{eq:functor-trace-factorization}: input dependence at the
$\mathcal I$-level, output uniqueness at the $P$-level, and
$\mathcal I$-good representatives.
\end{remark}

\section{Singular and unibranch boundaries}
\label{sec:functor-singular-boundary}

The general algebraic envelope and finite normal descent interfaces of
Part~IV retain $\mathrm{(SD)}_k$ of \Cref{warn:ch35-siu-source-gap}, unless
the required instance has an independent proof. Their scope includes
inseparable and nonreduced fibres, but they do not supply an analytic trace theorem
on a singular complex ambient space.

\begin{remark}[Unibranch trace remains open]
\label{rem:functor-unibranch-open}
For a reduced irreducible unibranch compact K\"ahler space $X$, the
general construction of $\Tr_Y$ would require:
\begin{enumerate}
\item quasi-equisingular regularization on the singular ambient space;
\item a multiplier-ideal transformation rule compatible with that
      regularization;
\item restriction estimates through a resolution independent of the
      chosen embedding; and
\item strong-topology completeness on the normalized target.
\end{enumerate}
The algebraic envelope property supplies none of these analytic inputs.
Accordingly, neither \Cref{thm:functor-split-trace} nor
\Cref{thm:functor-retraction-trace} is stated for a singular
ambient space.
\end{remark}

\begin{remark}[No general Dirichlet theorem]
\label{rem:functor-general-dirichlet-open}
For an arbitrary inclusion $Y\subset X$ and an arbitrary K\"ahler
class $\alpha$, extending a prescribed singularity while retaining
strict positivity in $\alpha$ remains unresolved. The product and
retraction theorems work because the cohomology class and the tangent
directions split. They are special cases, not reductions of the
general problem.
\end{remark}

\begin{remark}[Finite descent and analytic trace]
\label{ex:functor-finite-does-not-repair}
Let $\nu:X^\nu\to X$ be the normalization of a unibranch compact
complex space. The map on analytifications may have excellent
topological properties, and determinant norms may descend algebraic
metrics. Nevertheless, a trace built upstairs is not known to be
independent of all analytic choices downstairs without the
equisingular regularization package listed above.
\end{remark}

\section{Functoriality checklist}
\label{sec:functor-checklist}

\begin{proposition}\label{prop:functor-operation-ledger}
The valid operations and their inputs are:
\begin{enumerate}
\item proper bimeromorphic maps admit pullback of currents and
      singularity types;
\item finite pushforward for continuous semipositive metrics gives
      $d\,\operatorname{MA}(\phi)$;
\item finite locally free algebras admit the length-weighted
      determinant norm;
\item analytic trace requires $\nu(\varphi,Y)=0$ and produces a
      canonical $P$-class on $\widetilde Y$, represented by
      $\mathcal I$-good potentials and depending only on the ambient
      $\mathcal I$-class;
\item the split Dirichlet solution requires a product splitting or the
      transverse-positive retraction splitting.
\end{enumerate}
\par
\end{proposition}

\begin{proof}
The rows are respectively
\Cref{prop:functor-birational-pullback},
\Cref{prop:functor-finite-ma},
\Cref{prop:functor-determinant-norm},
\Cref{def:functor-analytic-trace}, and
\Crefrange{thm:functor-split-trace}{thm:functor-retraction-trace}.
Each required hypothesis appears in the proof at the point where its
omission would invalidate the construction.
\end{proof}

\section*{Source notes}
\label{sec:functor-source-notes}

Birational arguments use \Cref{chap:complex-prerequisites} and
\cite[Remark~1.7 and Proposition~1.12]{BEGZ10}. The arbitrary-singularity
Jacobian identity and continuous determinant norm are proved here; the norm
includes residue degrees and coalescing fibres. Flattening follows
\Cref{chap:bj-finite-descent}. The valuation-function base-change operator
is \cite[\S6.3, equation~(6.2)]{BHJ17}.

The trace interfaces are \Cref{thm:operations-trace-well-defined} and
\Cref{thm:operations-trace-p-descent}, using
\cite[Corollary~7.1.2 and \S8.1--8.2]{XiaBook}. The supplied book provides
$P$-class uniqueness and the general background interface. The accepted
article supplies analytic trace; its constant-approximation interface
\cite[Example~3.7]{DXTraceV2} requires rational analytic singularities.
The historical product audit was at the $\mathcal I$-level; its displayed
$P$-class follows from sourced uniqueness and is not a new production
receipt. Retraction requires the stated splitting and transverse
positivity. Unibranch trace and the general Dirichlet problem remain open.
